\section{Bounded nesting: the owner lemma and Theorem F}\label{sec:nesting}

The fourth paper's Theorem~B cuts the walk DAG of any coupling of home-based walks with a lexicographic label: one
digit per occupied group of scales, each counting legs measured from a fixed home~\cite[Theorem~4.20]{FourthPaper2026}.
Its hypothesis, linear leg area about fixed homes, fails when heads turn often, far from every home, on legs much shorter
than their distance from home. This section
replaces homes by owners. \Cref{lem:owner} shows that nested structures with distinct turning points force revisits, and
Theorem~F (\cref{thm:F}) proves depth reduction for every family of walks of bounded \emph{owner-load}
(\cref{def:owner-load}), whatever their homes, drift and scales. Its label is Theorem~B's band label run on \emph{virtual
heads}, one per head, side and nesting depth, with the band of a virtual head given by the height of its current family of
hops instead of a distance from home. \Cref{sec:nest-capacity} derives the consequences at walk level: TCP for walks of
bounded entry, and the confinement of depth-robustness to a set of cells of vanishing capacity.

\subsection{The owner lemma}\label{sec:nest-owner}

\begin{lemma}[Nesting with distinct owners costs visits]\label{lem:owner}
\formal{Owner.lean}{owner_entries, card_crossings_le}
Let $x$ be a walk, $\ell\ge1$, and let $(p_i,u_i)$, $0\le i<\ell$, be upper long hops of~$x$ that are strictly nested,
$p_i<p_{i+1}$ and $u_{i+1}<u_i$, with pairwise distinct owners. Then for every time $T'\ge u_0$, the cell of the innermost
hop $(p_{\ell-1},u_{\ell-1})$ has at least $(\ell+1)/2$ entries at times at most~$T'$. By the reflection $x\mapsto-x$, the
same holds for lower long hops.
\end{lemma}

\begin{outline}
The Lean statement concerns upper hops; the lower case follows by the reflection $x\mapsto-x$, which maps lower long hops
to upper ones with the same owners and does not change entries. In Lean a walk is a map $x\colon\N\to\Z$ with
$\abs{x(t+1)-x(t)}\le1$ (\lean{IsWalk}); a walk of length~$T$ is extended by staying at $x(T-1)$, which adds no entries
before time~$T$. Upper long hops and owners are the structures
\lean{IsUpperHop} and \lean{IsOwner}, and entries and crossings are \lean{entries} and \lean{crossings}. Let $c_i$ be the
cell of $(p_i,u_i)$ and $c=c_{\ell-1}$. Since $p_{i+1}\in(p_i,u_i)$, the cells increase inwards, $c_0<\dots<c_{\ell-1}=c$.
Every owner $o_i$ is a maximum of $x$ over $(p_i,u_i)\supseteq(p_{\ell-1},u_{\ell-1})$, where $x\ge c+1$, so
$x(o_i)\ge c+1$. For $i<\ell-1$, $o_i$ lies outside $[p_{i+1},u_{i+1}]$: at the endpoints $x=c_{i+1}\le c$, and if $o_i$
lay strictly inside $(p_{i+1},u_{i+1})$ it would also be the first maximum over that interval, that is $o_{i+1}$
(\lean{IsOwner.unique}), contrary to distinctness. So the walk crosses the edge between $c$ and $c+1$ at some step in
$(p_i,p_{i+1})$, between $o_i$ and $p_{i+1}$, or in $[u_{i+1},u_i)$, between $u_{i+1}$ and~$o_i$. The innermost hop
crosses it at the steps $p_{\ell-1}$ and $u_{\ell-1}-1$. These $\ell+1$ crossings occur at distinct steps, because the
intervals $(p_i,p_{i+1})\cup[u_{i+1},u_i)$ are pairwise disjoint and disjoint from $[p_{\ell-1},u_{\ell-1})$. Finally,
every crossing of the edge either enters~$c$ or leaves it, each maximal run of visits to~$c$ is entered once and left at
most once, and so the number of crossings is at most twice the number of entries (\lean{card_crossings_le}).
\end{outline}

\begin{corollary}[Bounded entry bounds the owner-load]\label{cor:owner-load}
\notformalized
Let $x_1,\dots,x_k$ be walks of length~$T$ and $H'$ a set of their long hops such that, for every hop of~$H'$ of head~$h$
at a cell~$c$, the walk $x_h$ enters~$c$ at most $V$ times. Then $H'$ has owner-load at most $2V-1$. In particular, if every
walk enters every cell at most $V$ times, the set of all long hops has owner-load at most $2V-1$.
\end{corollary}

\begin{informalproof}
Fix a time~$t$, a head~$h$ and a side, say upper. The hops of~$H'$ of head~$h$ on this side whose open intervals
contain~$t$ form a chain (\cref{sec:prelim-walks}). Choosing one hop for each of their $\ell$ distinct owners gives a strictly
nested chain of upper long hops with distinct owners, and by \cref{lem:owner} with $T'=T-1$ its innermost cell has at
least $(\ell+1)/2$ entries. This cell is the cell of a hop of~$H'$ of head~$h$, so $(\ell+1)/2\le V$. Lower hops are
symmetric.
\end{informalproof}

\subsection{Families and virtual heads}\label{sec:nest-families}

\begin{definition}[Families]\label{def:families}
\formal{Families.lean}{side, sw, fam, famP, famU, famd, famDepth, vcode, famData}
Let $H'$ be a set of long hops of walks $x_1,\dots,x_k$ of length~$T$. For a head~$h$, a side~$\sigma$ and a time~$o$,
the \emph{family} of $(h,\sigma,o)$ is the set of hops of~$H'$ of head~$h$ on side~$\sigma$ with owner~$o$, when it is
nonempty; every hop of~$H'$ lies in exactly one family. All hops of a family~$F$ contain its owner~$o_F$, so they form a
chain. The \emph{span} of~$F$ is its outermost hop $(P_F,U_F)$, at the cell~$c_F$; its \emph{height} is
$d_F=\abs{x_h(o_F)-c_F}\ge1$; its \emph{band} is $\lambda(d_F)$, where $\lambda(d)=\lfloor\log_2d\rfloor$, so
$2^{\lambda(d)}\le d<2^{\lambda(d)+1}$; and its \emph{depth} is the number of families of $(h,\sigma)$ whose span
contains the span of~$F$ as an interval, $F$ included. A \emph{virtual head} is a triple $v=(h,\sigma,i)$ such that some
family of $(h,\sigma)$ has depth~$i$; its families are those of $(h,\sigma)$ of depth~$i$. We write $N$ for the number of
virtual heads.
\end{definition}

\begin{lemma}[Structure of families]\label{lem:families}
\formal{Families.lean}{fam_cell_mem, fam_cell_injOn, card_fam_le, fam_long, famP_ne, fam_no_cross, famDepth_le, vcode_lt, vh_eq, famData_disj, famData_valid}
\formal{Clock.lean}{Fam.Valid.unique_span, Fam.Valid.sum_d_le, Fam.Valid.overlap}
Let $H'$ have owner-load at most~$\ell$.
\begin{enumerate}[label=(\alph*)]
\item The cells of the hops of a family~$F$ are distinct and lie between $c_F$, included, and $x_h(o_F)$, excluded, so
$\card F\le d_F$; moreover $U_F-P_F\ge2d_F$.
\item Distinct families of $(h,\sigma)$ have distinct spans, and their spans do not cross.
\item Every family has depth at most~$\ell$. Hence $N\le2k\ell$.
\item Two distinct families $F,F'$ of one virtual head satisfy $U_F\le P_{F'}$ or $U_{F'}\le P_F$.
\item For every virtual head~$v$ and time~$\tau$, at most one family~$F$ of~$v$ has $\tau\in(P_F,U_F]$, and
$\sum_{F\text{ of }v}d_F\le(T-1)/2$.
\item \emph{(Overlap.)} For every time~$\tau$ and integer $\theta\ge0$, at most $N(2^\theta-1)$ hops
$(p,u)\in H'$ with $p<\tau\le u$ lie in families of band less than~$\theta$.
\end{enumerate}
\end{lemma}

\begin{outline}
In Lean the walks are maps $x_h\colon\N\to\Z$ (\lean{IsWalk}), extended beyond~$T$ as in \cref{lem:owner}; the hops of~$H'$ are
triples $(h,p,u)$ such that $(p,u)$ is a long hop of~$x_h$ before~$T$ (\lean{IsHop}); and the owners may be any assignment of
times strictly inside the hops, the owners of \cref{sec:prelim-walks} being one (\lean{hopOwner_spec}). A hop is upper if the
walk steps up after its start (\lean{side}); seen from its side, that is, after the reflection $x\mapsto-x$ for lower hops
(\lean{sw}), it lies strictly above its cell (\lean{sw_above}), so all hops of a family are upper hops of one walk. Families,
spans, heights and depths are \lean{fam}, \lean{famP}, \lean{famU}, \lean{famd} and \lean{famDepth}; the span is the hop of
the family that starts first, and it also ends last (\lean{fam_nested}, \lean{fam_span}). With the heads numbered
$0,\dots,k-1$ and $\sigma=1$ for the upper side and $0$ for the lower one, a virtual head $(h,\sigma,i)$ is coded by the number
$(i-1)+(2h+\sigma)\ell<2k\ell$ (\lean{vcode}, \lean{vcode_lt}, \lean{vh_eq}). The formal family clock runs
on all $2k\ell$ codes, the unused ones being virtual heads without families, of band~$0$ at all times; so the formal~$N$ is
its bound~$2k\ell$. The data of all families form an instance \lean{famData} of abstract family data (\lean{Fam}), and
\lean{famData_valid} proves the hypotheses \lean{Fam.Valid} of the family clock: every hop lies in the span of its family with
the owner strictly inside, a hop that ends at the end of its span is the span, (d), (a), and spans end before~$T$. Parts (e)
and~(f) are proved from these hypotheses alone (\lean{Fam.Valid.unique_span}, \lean{Fam.Valid.sum_d_le},
\lean{Fam.Valid.overlap}).

(a) (\lean{fam_cell_mem}, \lean{fam_cell_injOn}, \lean{card_fam_le}, \lean{fam_long}.) Let $F$ be upper; the lower case is the same seen from its side. If $(p,u)\supsetneq(p',u')$ are hops of~$F$, then $p'\in(p,u)$, so
$x_h(p')>x_h(p)$: the cells increase strictly inwards. The innermost hop contains~$o_F$, where the walk is above its cell,
so all cells lie in $[c_F,x_h(o_F)-1]$, and $\card F\le d_F$. The walk moves by at most one cell per step and goes from
$c_F$ at $P_F$ to $x_h(o_F)$ at~$o_F$ and back to $c_F$ at~$U_F$, so $o_F-P_F\ge d_F$ and $U_F-o_F\ge d_F$.

(b) (\lean{famP_ne}, \lean{fam_no_cross}.) The span of~$F$ is a hop of~$F$, with owner~$o_F$; equal spans give equal owners
and hence equal families. Since a hop is determined by its start, distinct spans even start at distinct times. Spans are
long hops of one head on one side, which do not cross.

(c) (\lean{famDepth_le}.) Let $t\in(P_F,U_F)$, a nonempty interval by~(a). Every family~$F'$ whose span contains the span of~$F$ has $t$ in the
open interval of its span, a hop of~$H'$ of head~$h$ on side~$\sigma$ with owner~$o_{F'}$, and distinct such families have
distinct owners. So there are at most $\ell$ of them. A virtual head is then a triple $(h,\sigma,i)$ with $i\le\ell$.

(d) (\lean{famData_disj}, \lean{disj_aux}.) By~(b) we may assume $P_F<P_{F'}$. Suppose that $P_{F'}<U_F$. Since the spans do
not cross, $U_{F'}\le U_F$, and $U_{F'}\ne U_F$, since a hop is also determined by its end; so
$(P_{F'},U_{F'})\subsetneq(P_F,U_F)$. Every family whose span contains the span of~$F$ also contains the span of~$F'$, and
$F'$ contains its own span but not that of~$F$; so $F'$ would have larger depth than~$F$ (\lean{famDepth_lt}). Hence
$U_F\le P_{F'}$.

(e) (\lean{Fam.Valid.unique_span}, \lean{Fam.Valid.sum_d_le}.) The first claim is~(d). The half-open spans $(P_F,U_F]$ of the families of~$v$ are therefore disjoint subsets of
$\{1,\dots,T-1\}$, and by~(a) each has at least $2d_F$ elements.

(f) (\lean{Fam.Valid.overlap}.) A hop $(p,u)$ of a family~$F$ of a virtual head~$v$ with $p<\tau\le u$ satisfies $P_F\le p<\tau\le u\le U_F$. By~(e)
at most one family of~$v$ qualifies, and if its band is less than~$\theta$, then by~(a) it has at most
$d_F<2^{\lambda(d_F)+1}\le2^\theta$ hops. Summing over the $N$ virtual heads gives the bound.
\end{outline}

\begin{definition}[The family clock]\label{def:clock}
\formal{Clock.lean}{lam, grp, flr, Fam.band, Fam.group, Fam.occ, Fam.rk, Fam.cnt, Fam.digit, Fam.clock}
Let $H'$ have owner-load at most~$\ell$ and $N\ge1$ virtual heads, numbered $1,\dots,N$.
\begin{itemize}[leftmargin=2em]
\item The \emph{band} of a virtual head~$v$ at time~$t$ is $b_v(t)=\lambda(d_F)$ for the family~$F$ of~$v$ with the
largest $P_F\le t$, and $b_v(t)=0$ if no family of~$v$ starts at or before~$t$. It changes only at the start times
$P_F\ge1$ of the families of~$v$, the \emph{jumps} of~$v$.
\item The \emph{events} are the owners $o_F$ of all families of all virtual heads, one per family, with \emph{scale}
$\lambda(d_F)$.
\item Fix integers $a\ge1$ and $K\ge2$, an \emph{offset} $r\in\{0,\dots,a-1\}$ and $\omega\in\{0,\dots,K-1\}^N$. The
\emph{group} of~$v$ at~$t$ is $G_v(t)=\lfloor(b_v(t)+r)/a\rfloor$, and the \emph{floor} of a group $G\ge0$ is
$\theta(G)=\max(aG-r,0)$; a band $b$ lies in group~$G$ exactly when $\theta(G)\le b<\theta(G+1)$, and the floors at
least~$1$ are the integers congruent to $-r$ modulo~$a$.
\item The \emph{occupied groups} at~$t$ are the distinct values among $G_1(t),\dots,G_N(t)$, listed from the top as
$G_1(t)>\dots>G_{m(t)}(t)$, with floors $\theta_i(t)=\theta(G_i(t))$; here $m(t)\le N$.
\item The \emph{counts} are $n_i(t)=\#\{\text{events at times at most }t\text{ of scale at least }\theta_i(t)\}$.
\item The \emph{family clock} is the label
\[
  \Psi(t)=\sum_{i=1}^{m(t)}\bigl((n_i(t)+\omega_i)\bmod K\bigr)\,K^{N-i}\;\in\;[0,K^N).
\]
\end{itemize}
\end{definition}

So $\Psi(t)$ is the integer whose base-$K$ digits, most significant first, are $\psi_i(t)=(n_i(t)+\omega_i)\bmod K$ for
$i\le m(t)$ and $0$ for $m(t)<i\le N$, and labels compare as these digit strings compare lexicographically. Two bands
$b<b'$ lie in different groups exactly when some floor $\theta\ge1$ satisfies $b<\theta\le b'$; for a
uniformly random offset each integer $\theta\ge1$ is a floor with probability $1/a$. This is the band label of the fourth
paper's Theorem~B~\cite[Definition~4.19]{FourthPaper2026}, with virtual heads in place of heads, family bands in place of
distances from home, and owners in place of leg starts.

\subsection{Theorem F}\label{sec:nest-F}

\begin{theorem}[Theorem F: depth reduction at bounded owner-load]\label{thm:F}
\formal{TheoremF.lean}{theoremF_owner, theoremF, DF}
Let $k,\ell\ge1$, $\eps\in(0,1]$ and $T\ge1$, let $x_1,\dots,x_k$ be walks of length~$T$, coupled arbitrarily, and let $H'$
be a set of their long hops of owner-load at most~$\ell$. Then the DAG $G(H')$ has a cut of at most $\eps T$ edges and depth
less than
\[
  \DF(k,\ell,\eps)=2^{\,2k\ell\left(\lceil64k^2\ell^2/\eps\rceil+5+\lceil\log_2(k^2\ell^2/\eps)\rceil\right)}
  =2^{O(k^3\ell^3/\eps)}.
\]
The bound does not depend on~$T$, on the coupling, on homes or on the scales of the walks, and pauses are allowed.
\end{theorem}

\begin{outline}
The theorem for the owners of \cref{sec:prelim-walks} (\lean{hopOwner}) is \lean{theoremF_owner}. It is an instance of
\lean{theoremF}, which holds for every assignment of owners strictly inside the hops and for every $\eps>0$. Owner-load is
\lean{OwnerLoadLE}, and the hop edges of $G(H')$ are the pairs of \lean{hopPairs}. The cut consists of the state edges
$t\to t-1$ at the times of a subset of $[1,T)$ and of a set of hop pairs, and the depth bound concerns every path
$f(0)\to\dots\to f(n)$ that avoids the cut, as in \lean{strip_cut_depth}. The family clock is defined for abstract family data
(\lean{Fam.clock}) and applied to \lean{famData}, whose hypotheses \lean{Fam.Valid} hold by \cref{lem:families}, with
$N=2k\ell$ as explained there; digit positions and the coordinates of~$\omega$ are numbered from~$0$, so that position $i-1$
carries the digit~$\psi_i$. The expectation is formalized as a count: summed over the $a$ offsets and the $K^N$ choices
of~$\omega$, the violated hops and the descents number at most $K^N\cdot4N^2T+aK^{N-1}2^{a+1}N^2T$
(\lean{Fam.Valid.cost_sum_le}), which is $aK^N$ times the right side of \eqref{eq:F-expect}, and some $(r,\omega)$ does at most
as well as the average (\lean{Fam.Valid.exists_cheap}). The bounds below are lemmas of that count.

Let $r$ and~$\omega$ be uniformly random and independent, and let $F_\Psi$ be the cut of the family clock in $G(H')$
(violated hop edges and state edges at descents). By \cref{lem:labels} it has depth at most $K^N-1$; in Lean, for a label
with $N$ digits in $[0,K)$ compared lexicographically (\lean{Fam.depth_le_of_lex}, \lean{Fam.digit_lt}). We show
\begin{equation}\label{eq:F-expect}
  \mathbb E\card{F_\Psi}\;\le\;\frac{4N^2T}a+\frac{2^{a+1}N^2T}K.
\end{equation}
We bound the number of violated hops of~$H'$, counted per hop, which is at least the number of violated hop edges, and the
number of descents. Throughout, $\eta=(p,u)$ is a hop of~$H'$ of the virtual head~$w$, in the family~$F$, and
$G_\eta=\lfloor(\lambda(d_F)+r)/a\rfloor$ is the group of its band.

\emph{Consistency} (\lean{Fam.Valid.band_of_mem}, \lean{Fam.Valid.band_at_end}). We have $G_w(p)=G_\eta$, and $G_w(u)=G_\eta$
unless $\eta$ is the span of~$F$ and the next family $F'$ of~$w$ starts at $P_{F'}=U_F$. Indeed, $P_F\le p<u\le U_F$, and
by \cref{lem:families}(d) no family of~$w$ starts in $(P_F,U_F)$, so $b_w=\lambda(d_F)$ on $[P_F,U_F)$; the hops of~$F$
other than its span have $u<U_F$, by strict nesting. Call $\eta$ \emph{inconsistent} if $G_w(u)\ne G_\eta$.

\emph{The jumps that change a group} (\lean{Fam.gjumps}, \lean{Fam.Valid.sum_gjumps}). A jump of~$v$ from band~$b$ to
band~$b'$ changes the group of~$v$ with probability at most $\abs{b-b'}/a\le(b+b')/a$ (\lean{card_grp_ne_le}). At the jump that starts a
family~$F''$, the new band is $\lambda(d_{F''})$ and the old band is that of the previous family of~$v$, or the initial
band~$0$; each family occurs at most once as a new band and at most once as an old band (\lean{Fam.Valid.sum_old_le}). So,
since $\lambda(d)<d$, $\sum_{\text{jumps of }v}(b+b')\le2\sum_F\lambda(d_F)\le T$ (\lean{Fam.Valid.sum_jumps_lam}), and
the jumps that change a group number at most $NT/a$ in expectation.

\emph{($\alpha$) Inconsistent hops} (\lean{Fam.incons}, \lean{Fam.Valid.card_incons_le}). We count them all as violated. By
consistency, an inconsistent hop is the span of its family~$F$, and the next family $F'$ of~$w$ starts at $U_F$ with a jump
that changes the group of~$w$; distinct inconsistent hops give distinct such jumps. So their expected number is at most
$NT/a$.

\emph{($\beta$) Consistent hops whose higher groups differ} (\lean{Fam.higher}, \lean{Fam.Valid.exists_charge},
\lean{Fam.Valid.card_higher_le}, \lean{Fam.Valid.sum_higher_le}). Let $\mathcal U(t)$ be the set of occupied groups above
$G_\eta$ at time~$t$, and suppose $\eta$ is consistent and $\mathcal U(p)\ne\mathcal U(u)$; we count it as violated. Some group
$G'>G_\eta$ is occupied at exactly one of the times $p,u$, by a virtual head~$v$ that is not in~$G'$ at the other time;
$v\ne w$, since $w$ is in $G_\eta$ at both times. So $G_v(p)\ne G_v(u)$ and $\max(G_v(p),G_v(u))>G_\eta$. The group of~$v$
changes only at its jumps, and so some jump $\tau\in(p,u]$ of~$v$ moves it between two groups whose larger one, $G^\circ$,
exceeds~$G_\eta$ (\lean{exists_change}: going back from~$u$, the first step at which the group of~$v$ changes between two
groups the larger of which exceeds~$G_\eta$). Let the jump move the band of~$v$ from~$b$ to~$b'$. The floor
$\theta^\circ=\theta(G^\circ)$ satisfies $\min(b,b')<\theta^\circ\le\max(b,b')$, and since
$G^\circ\ge G_\eta+1$, $\lambda(d_F)<\theta(G_\eta+1)\le\theta^\circ$. So $\eta$ spans~$\tau$ and its family has band
less than~$\theta^\circ$. Charge $\eta$ to the jump $(v,\tau)$. By the overlap bound of \cref{lem:families}, at most
$N(2^{\theta^\circ}-1)$ hops are charged to one jump. Given~$r$, $\theta^\circ$ is the largest floor in
$(\min(b,b'),\max(b,b')]$, and a jump that changes no group receives no charge. Over the $a$ offsets at which the jump
changes a group, these floors are distinct integers in $(\min(b,b'),\max(b,b')]$, since the floors at least~$1$ for the
offset~$r$ are the integers congruent to $-r$ modulo~$a$ (\lean{sum_floor_pow_le}); so the expected charge of a jump is at
most $(N/a)2^{M+1}$, where $M=\max(b,b')$ and $2^M\le2^b+2^{b'}$. As above, each family occurs at most once as a new band
and at most once as an old band, so $\sum_{\text{jumps of }v}2^M\le1+2\sum_F2^{\lambda(d_F)}\le1+2\sum_Fd_F\le T$
(\lean{Fam.Valid.sum_jumps_pow}). The expected number of such hops is therefore at most $(2N/a)\sum_vT=2N^2T/a$.

\emph{($\gamma$) The remaining hops} (\lean{Fam.calm}, \lean{Fam.Valid.card_viol_calm_le}, \lean{Fam.Valid.sum_calm_growth_le},
\lean{Fam.Valid.sum_viol_calm_le}). Now $\eta$ is consistent and $\mathcal U(p)=\mathcal U(u)$. Then $G_\eta$ has the same
position $j=\card{\mathcal U(p)}+1$ in the lists of occupied groups at $p$ and at~$u$, and the groups above it, hence
their floors $\theta_i$, $i\le j$, agree. Put $\Delta_i=n_i(u)-n_i(p)$ for $i\le j$; then $\Delta_i\ge0$, because the
floor is the same at both times. Moreover $\Delta_j\ge1$: the owner $o_F\in(p,u)$ is an event of scale
$\lambda(d_F)\ge\theta(G_\eta)=\theta_j$. Say that digit~$i$ \emph{wraps} if
$\lfloor(n_i(p)+\omega_i)/K\rfloor\ne\lfloor(n_i(u)+\omega_i)/K\rfloor$. If no digit $i\le j$ wraps, then
$\psi_i(u)=\psi_i(p)+\Delta_i$ for $i\le j$, the first position $i\le j$ with $\Delta_i>0$ carries a larger digit at~$u$,
all earlier digits agree, and $\Psi(u)>\Psi(p)$ (\lean{Fam.clock_lt_of_no_wrap}). Given~$r$, digit~$i$ wraps with probability
at most $\Delta_i/K$ (\lean{Fam.card_wrap_le}, \lean{Fam.card_wrap_union}), so $\eta$ is violated with probability at most
$\sum_{i\le j}\Delta_i/K$.

We bound $\sum_\eta\sum_{i\le j(\eta)}\Delta_i(\eta)$ over these hops, for every~$r$, by exchanging the order of
summation. $\Delta_i(\eta)$ is the number of events $o'\in(p,u]$ of scale at least~$\theta_i$. Fix an event~$o'$ of
scale~$s$. A pair $(\eta,i)$ counts~$o'$ only if $\eta$ spans~$o'$ and $\theta(G_\eta)\le\theta_i\le s$; the floors
$\theta_i$, $i\le j(\eta)$, are distinct, since $\theta$ is strictly increasing on groups. If $\theta(G_\eta)\le\varphi$,
then $\lambda(d_{F(\eta)})<\theta(G_\eta+1)\le\theta(G_\eta)+a\le\varphi+a$. So, by the overlap bound with $\theta=\varphi+a$,
the pairs
counting~$o'$ number at most
\[
  \sum_{\varphi}N\bigl(2^{\varphi+a}-1\bigr)\;<\;N2^{a+1}2^s,
\]
the sum over the floors $\varphi\le s$, which are distinct integers in $[0,s]$. Each family is the family of one event, and
$\sum_{o'}2^{s(o')}=\sum_v\sum_F2^{\lambda(d_F)}\le\sum_v\sum_Fd_F\le NT/2$. Hence these hops contribute at most
$N2^{a+1}\cdot NT/(2K)=2^aN^2T/K$ in expectation.

\emph{Descents} (\lean{Fam.still}, \lean{Fam.card_desc_still_le}, \lean{Fam.sum_desc_still_le}, \lean{Fam.gchange},
\lean{Fam.Valid.card_gchange_le}). Consider the state edge $t\to t-1$. If no virtual head changes group at step $t-1\to t$, the
occupied groups and their floors are the same at $t-1$ and~$t$, and $n_i(t)\ge n_i(t-1)$ for every~$i$; as in~($\gamma$), a
descent needs a wrap (\lean{Fam.clock_le_of_no_wrap}), which given~$r$ has probability at most $\sum_i(n_i(t)-n_i(t-1))/K$.
Summed over~$t$ this is at most $N$ times the number of events divided by~$K$, and there are at most
$\sum_v\sum_Fd_F\le NT/2$ events; so these descents number at most $N^2T/(2K)$ in expectation. Otherwise some virtual head
has a jump at~$t$ that changes its group, and we count the edge as cut; these edges number at most $NT/a$ in expectation,
by the bound on the jumps that change a group.

\emph{Total} (\lean{Fam.violated_subset}, \lean{Fam.descents_subset}, \lean{Fam.Valid.cost_sum_le}). Adding the four bounds,
$\mathbb E\card{F_\Psi}\le(2NT+2N^2T)/a+(2^a+1/2)N^2T/K$, which gives \eqref{eq:F-expect}. (The formal count bounds the two
terms with~$K$ together: by $N2^{a+1}\sum_Fd_F/K$ and by $N/K$ per family, at most $2^{a+1}N^2T/K$ in all.)

\emph{Parameters} (\lean{theoremF}, \lean{cost_le_eps}, \lean{ceil_le_two_pow_clog}). By \cref{lem:families}(c), $N\le2k\ell$,
so the right side of \eqref{eq:F-expect} is at most $16k^2\ell^2T/a+2^{a+3}k^2\ell^2T/K$. Take
$a=\lceil64k^2\ell^2/\eps\rceil$ and $K=2^{a+4}\lceil k^2\ell^2/\eps\rceil$.
Then the first term is at most $\eps T/4$ and the second at most $\eps T/2$, and some choice of $(r,\omega)$ gives a cut of
at most $\eps T$ edges. Its depth is at most $K^N-1<K^{2k\ell}$, and since $\lceil x\rceil\le2^{\lceil\log_2x\rceil}$
for $x>0$,
\[
  K^{2k\ell}\le2^{2k\ell\left(\lceil64k^2\ell^2/\eps\rceil+4+\lceil\log_2(k^2\ell^2/\eps)\rceil\right)}<\DF(k,\ell,\eps).
\]
In \lean{DF} the ceiling of the logarithm is taken in~$\N$, which agrees with the integer ceiling for $\eps\le1$, since
$k^2\ell^2/\eps\ge1$. Pauses are state edges, and the argument uses only that hops are long and that owners lie strictly
inside their hops.
\end{outline}

\paragraph{Why it matters.} Theorem~F proves the walk form of the structural hypothesis of Theorem~D for every family of
bounded owner-load, for every coupling, and with no reference to positions: neither homes nor leg area enter. Together
with the owner lemma it shows that a depth-robust family of walks, if one exists, must contain heads with unboundedly
many simultaneously open nested structures with distinct turning points, hence, by \cref{lem:owner}, cells entered
unboundedly often. This is the structure of the deep-nesting pair (\cref{sec:prelim-dnp}), which Theorem~F does not
cut at a depth independent of its nesting depth. Theorem~F is consistent with the fourth paper's nesting barrier for band
rules~\cite[Remarks~5.6 and~5.7]{FourthPaper2026}, which concerns rules whose load does not grow with the nesting depth;
the load of the family clock is $N\le2k\ell$.

\subsection{Bounded entry and the capacity of depth}\label{sec:nest-capacity}

\begin{corollary}[Bounded entry implies TCP]\label{cor:entry-tcp}
\notformalized
Let $V\ge1$. TCP holds for every family of $k$-tuples of walks in which every walk enters every cell at most $V$ times,
in particular for walks of multiplicity at most~$V$, with $D(\eps)=\DF(k,2V-1,\eps)$ for $\eps\in(0,1]$.
\end{corollary}

\begin{informalproof}
By \cref{cor:owner-load} the set of all long hops has owner-load at most $2V-1$, and $G(H')$ is the walk DAG; apply
\cref{thm:F}.
\end{informalproof}

The corollary is incomparable with Theorem~B: home-based sweeps $0\to n\to0$, repeated, enter cells near home unboundedly
often, while walks made of V's of varying sizes that drift away from every home enter every cell at most three times
but have superlinear leg area about every cell.

\begin{corollary}[The capacity of depth]\label{cor:capacity}
\notformalized
Let $x_1,\dots,x_k$ be walks of length~$T$, let $\eps\in(0,1]$ and $V=\lceil1/\eps\rceil$, and let $R$ be the set of pairs
$(h,c)$ such that $x_h$ enters~$c$ more than $V$ times. Then $\card R\le kT/(V+1)<k\eps T$; at every time~$\tau$ at most
$\card R$ hops at pairs of~$R$ span~$\tau$; and $G-R$ has a cut of at most $\eps T$ edges and depth less than
$\DF(k,2\lceil1/\eps\rceil-1,\eps)=2^{O(k^3/\eps^4)}$.
\end{corollary}

\begin{informalproof}
Each walk has at most $T$ entries in all, which bounds $\card R$. By (P1), at most one hop at each pair spans a
given time. The DAG $G-R$ is $G(H')$ for the set $H'$ of long hops at pairs entered at most $V$ times, which by
\cref{cor:owner-load} has owner-load at most $2V-1\le3/\eps$; apply \cref{thm:F}.
\end{informalproof}

\paragraph{Why it matters.} Every schedule of linear length splits, with no property of transposition used, into a
\emph{shallow} part, the hops at boundedly entered cells, which a cut of $\eps T$ edges reduces to bounded depth, and a
\emph{scratch}~$R$ of fewer than $k\eps T$ cells, which carries fewer than $k\eps T$ symbols across every time cut.
Paths that avoid the cut may still alternate between at most $\DF$ shallow hops and hops through the scratch. So
depth-robustness of a linear-time schedule can only live in a set of cells of vanishing capacity that its heads visit
over and over. On the deep-nesting family the trade-off between $\card R/T$ and the depth is of order $1/\log D$, and no
better under the walk-DAG form of Conjecture~A$'$ (\cref{prop:dnp-capacity}). The next proposition shows that the
alternation cannot be forbidden at no cost: for walks without pauses, the statement $\mathrm{RF}(k+1)$ that forbids it
implies $\TCP(k)$ in genie form, while $\TCP(k+1)$ in genie form implies $\mathrm{RF}(k+1)$ with $R=\emptyset$.

\begin{proposition}[Relays must be allowed]\label{prop:relays}
\notformalized
Restrict to walks without pauses, and let $\mathrm{RF}(k)$ be the statement: for every $\eps>0$ there are $D(\eps)$ and
$\delta(\eps)$, with $\delta(\eps)\to0$ as $\eps\to0$, such that for all $k$ walks of length~$T$ there are a set~$R$ of pairs
with $\card R\le\delta(\eps)T$ and a set $\Gamma\subseteq[0,T)$ with $\card\Gamma\le\eps T$ such that, for every path
$t_0\to\dots\to t_n$ of the walk DAG with $t_1,\dots,t_n\notin\Gamma$ and every $i$ with $D(\eps)<i\le n$, if
$t_{i-1}\to t_i$ is a hop edge then every hop that it carries lies at a pair of~$R$. Then $\mathrm{RF}(k+1)$ implies
$\TCP(k)$ in genie form: for every $\eps'>0$ there is~$D'$ such that every $k$ walks of length~$T$ have a set
$\Gamma'\subseteq[0,T)$ with $\card{\Gamma'}\le\eps'T$ such that every path whose vertices after the first avoid~$\Gamma'$
has at most $D'$ edges.
\end{proposition}

\begin{informalproof}
For $\eps'\ge1$ take $\Gamma'=[0,T)$. Let $0<\eps'<1$, let $M$ be the least odd integer with $M\ge4/\eps'$, and choose
$\eps\le\eps'/8$ with
$(M-1)\delta(\eps)\le\eps'/8$. Let $x_1,\dots,x_k$ be walks of length $T\ge16/\eps'$. Define an extra walk~$z$ on
$[0,T)$ blockwise: for $t\in[jM,(j+1)M)$, $z(t)=j+((t-jM)\bmod2)$. Inside block~$j$ it alternates between $j$ and $j+1$,
ending at~$j$ because $M$ is odd, and the next block starts at $j+1$; so $z$ has no pauses. Cell~$j$ is visited at most
$M$ times, $z$ has at most $T/M+2$ first visits, and every time that is not a first visit of~$z$ is the later endpoint of a
long hop of~$z$.

Apply $\mathrm{RF}(k+1)$ to $(x_1,\dots,x_k,z)$ at~$\eps$, obtaining $R$ and~$\Gamma$. Let $\Delta$ be the set of times~$s$
at which some path of length greater than $D(\eps)$, with all vertices after the first outside~$\Gamma$, ends. For such a
path ending at~$s$, if $s$ were not a first visit of~$z$, the pair $(z,z(s))$ were not in~$R$ and $\prev_z(s)$ were not
in~$\Gamma$, then appending the hop edge $s\to\prev_z(s)$ would give a path that violates $\mathrm{RF}$ at its last edge.
So every $s\in\Delta$ is a first visit of~$z$ (at most $T/M+2$ times), or the later endpoint of a hop of~$z$ at a pair
of~$R$ (at most $(M-1)\card R\le(M-1)\delta(\eps)T$ times, since each pair of~$z$ carries at most $M-1$ hops), or has
$\prev_z(s)\in\Gamma$ (at most $\card\Gamma\le\eps T$ times, since $\prev_z$ is injective).

Put $\Gamma'=\Gamma\cup\Delta$. A path of length greater than $D(\eps)$ whose vertices after the first avoid~$\Gamma'$
would avoid~$\Gamma$, so its last vertex would lie in~$\Delta\subseteq\Gamma'$, a contradiction. Hence such paths have at
most $D(\eps)$ edges in the walk DAG of $(x_1,\dots,x_k,z)$, and therefore in that of $(x_1,\dots,x_k)$, which has the same
vertices and fewer edges. Finally
$\card{\Gamma'}\le2\eps T+T/M+2+(M-1)\delta(\eps)T\le(\eps'/4+\eps'/4+\eps'/8)T+2\le\eps'T$, using $T\ge16/\eps'$. For
$T<16/\eps'$ every path has fewer than $16/\eps'$ edges, so $D'=\max(D(\eps),\lceil16/\eps'\rceil)$ works.
\end{informalproof}

The conversion is the reason the scratch of \cref{cor:capacity} must be allowed to relay: a busy extra head of bounded
multiplicity and space about $T/M$ turns any relay-free certificate into a genuine genie for the other heads.

\begin{remark}[Depth-robust schedules carry capacity]\label{rem:abp}
Proof sketch only. Alwen, Blocki and Pietrzak showed that a DAG that keeps a path with $d$ vertices after the deletion of
any $e$ vertices has cumulative pebbling cost greater than about $ed$~\cite{ABP2017}; graphs of this kind go back to Erdős,
Graham and Szemerédi~\cite{EGS1975}. Pebbling the walk DAG of $k$ walks in time order keeps, at step~$t$, pebbles on the
current and the previous step and on the earlier endpoints of the hops that span~$t$, at most $\Live(t)+2$ pebbles, where
$\Live(t)$ is the number of hops spanning~$t$. Hence, if a family is depth-robust at $(\eps,D)$ in genie form, then
$\sum_t\Live(t)$ is at least about $\eps TD-2T$: the mean time-cut capacity of a depth-robust schedule grows linearly with
the depth that it forces. \Cref{cor:capacity} bounds the capacity of the deep part from above; on the deep-nesting pairs
both bounds are compatible, with capacity about $T/(m+1)$ and, under Conjecture~A$'$, depth exponential in~$m$.
\end{remark}
